(1) 两个无相互作用、无自旋的全同粒子束缚在宽度 \(L\) 的一维无限深阱内，单粒子能级为 \(\epsilon_n=n^2\epsilon_1\)，\(\epsilon_1=\pi^2\hbar^2/(2mL^2)\)。分别求两个玻色子和两个费米子的最低总能量及占据。
(2) 若一粒子在 \(n=1\)，另一粒子在 \(n=2\)，写出归一化的玻色和费米双粒子空间波函数。
(3) 比较两种态在同一空间点 \(x_1=x_2=x\) 的联合概率密度，解释无自旋费米子为何不能同时占据同一个单粒子轨道。
(4) 定义单粒子配分函数 \(Z_1(\beta)=\sum_{n\ge1}e^{-\beta\epsilon_n}\)，推导两全同玻色子和两全同费米子的正则配分函数。
